\section{Why Retain Unfinished Exploration?}\label{sec:explanation}
\subsection{Local exploration needs room to develop}\label{sec:ablation}
A full-cohort factorial experiment varies scout admission and distance stopping on the 7,999-query, three-graph cohort. It is a retrospective mechanism extension after that cohort was exposed, with no new formal timing. Every policy uses the same $p+6{,}400$ actual distance work per event; each scout retains its 400-evaluation cap, including upper descent. L denotes strict local admission. GLOBAL stops against the main collector's threshold, LOCAL against the scout's local threshold, and NONE removes distance-based early stopping while retaining the cap and eligibility rules.

\begin{table}[t]
\centering
\caption{Scout stopping rules on 7,999 queries and three graphs (23,997 events). All policies match actual work per event. Bottom DC is mean new bottom-layer scoring in the scout phase; it excludes upper descent.}
\label{tab:stop}
\small
\begin{tabular}{@{}lrrrr@{}}
\toprule
Policy & Recall (\%) & Severe & Zero & Bottom DC \\
\midrule
C & 99.02654 & 65 & 25 & 0.00 \\
L\_GLOBAL & 99.09739 & 46 & 15 & 0.36 \\
L\_LOCAL (F) & 99.18240 & 26 & 10 & 1098.68 \\
L\_NONE & 99.17031 & 24 & 10 & 1342.70 \\
\bottomrule
\end{tabular}
\end{table}


Table~\ref{tab:stop} shows why immediately applying a mature global threshold can curtail a new route. L\_GLOBAL spends only 0.36 new bottom-layer distances per event across its scouts, compared with 1,098.68 for F. Their mean upper-descent costs are similar (254.86 and 253.86). F raises recall and reduces severe events from 46 to 26, with 22 rescues and two harms. The risk reduction is 0.08334 percentage points, with a descriptive query-cluster 95\% interval of [0.03750, 0.13335]. Unused scout work goes to final continuation, so the intervention tests a complete allocation policy, not stopping at equal scout expenditure.

L\_NONE has 24 severe events, two fewer than F, but slightly lower mean recall. Thus the local stopping rule is not uniquely effective or best on every metric. Allowing all discovered neighbors into the scout queue (A\_LOCAL) instead of strict local admission produces the same scoring paths and returned sets as F on all 23,997 events. A constructed distance-tie example distinguishes these rules; their observed agreement is not a theorem. The full factorial results in Appendix~\ref{app:factorial} prevent attributing the benefit to an admission filter unsupported by the ablation.

\subsection{Answers and continuation opportunities differ}
We use selected execution interventions to explain the design after establishing its overall effect. An earlier study selected 55 rescued and five harmed T2I events. For 22 of the 55 rescues, combining \emph{all} candidates scored by a complete C search with all candidates scored by scouts fails to reproduce the rescue. This is stronger than combining returned top-ten lists: it retains every scored candidate, yet lacks later discoveries made during F's continuation.

The strengthened union control uses 13,893.08 evaluations on average, exceeding F's 12,800. It tests candidate sufficiency, not equal-work performance. A separate intervention retains executable continuation states and compares six schedules at exactly 12,800 evaluations per event. Splitting remaining continuation 25\% or 50\% to the C state and running C before F preserves all 22 later rescues. Running F before C preserves 21 and 18, respectively. Thus feedback can enable later discoveries, while execution order still determines which opportunities are realized. These diagnostic schedules differ from F's production merge and do not prove that every operation is necessary.

The selected events do not estimate mechanism frequency in the full population. Search order, shared scores, visitation, and continuation state interact. The inference is that preserving an opportunity for subsequent discovery can matter, not that a uniquely identified queue or edge explains every improvement.

\subsection{Scouts spend opportunities on the original route}
Retaining more state does not refund the cost of exploration. We examine the recorded C/F prefixes while preserving the full original scout prefix and charging overlapping distance work once. Restoring five hits requires at least 13,859, 13,298, and 13,268 evaluations in three of the five harmed events. Each exceeds the original 12,800 allowance. Reallocating work among those recorded prefixes cannot eliminate these harms within that allowance.

This is a restricted-family result, not a lower bound for arbitrary graph searches. A different entrance, expansion order, scout requirement, or larger allowance may change the outcome. Its design implication is concrete: alternative exploration must be judged against the useful main-route work it displaces. The LAION and WebVid outcomes demonstrate why keeping more scout state or adding a selection stage is insufficient by itself.
